\chapter{Engenharia de Software}
\label{cap:eng-software}

\begin{edital}
Engenharia de requisitos (classificação, processo, técnicas de elicitação);
testes (unitários, integração, ágeis, usabilidade, automatizados, TDD, ciclo
de vida, RPA); metodologias ágeis (Scrum, Kanban, XP); padrões de
desenvolvimento e reuso; codificação; Ponto de Função e Story Points; DevOps;
design de software.
\end{edital}
\peso{MUITO ALTO --- foi o maior bloco da Dataprev 2024 (${\sim}$10 questões).
É metade do ``eixo duplo'' da FGV em TI (a outra metade é Banco de Dados).}

\section{Ciclo de vida e modelos de processo}

\begin{conceito}[Os modelos clássicos]
\begin{itemize}
\item \textbf{Cascata (waterfall):} fases sequenciais estritas, sem
sobreposição, requisitos congelados no início. Adequado quando os requisitos
são \textbf{estáveis e bem compreendidos}. Rígido a mudanças.
\item \textbf{Incremental:} entrega em incrementos que somam funcionalidade
ao longo do tempo.
\item \textbf{Iterativo:} refina o mesmo produto em ciclos sucessivos.
\item \textbf{Espiral (Boehm):} iterativo + \textbf{análise de risco} a cada
volta da espiral.
\item \textbf{Prototipação:} constrói um protótipo (às vezes descartável)
para validar requisitos com o cliente.
\item \textbf{Ágil:} iterativo-incremental, com entregas curtas e adaptação
contínua a mudanças.
\end{itemize}
\end{conceito}

\begin{pegadinha}
O enunciado descreve ``fases sequenciais, requisitos congelados'' $\to$ é
\textbf{cascata}, nunca incremental nem ágil. Esse é o distrator clássico de
abertura do bloco.
\end{pegadinha}

\section{Engenharia de requisitos}

O par que mais cai no bloco inteiro (Dataprev, MPU e TJ-RJ todos cobraram):

\begin{center}
\begin{tabular}{@{}p{2.6cm}p{5.4cm}p{5cm}@{}}
\toprule
\textbf{Tipo} & \textbf{O que descreve} & \textbf{Exemplo} \\
\midrule
\textbf{Funcional (RF)} & O QUE o sistema faz (uma função, um comportamento) & ``emitir saldo da conta'' \\
\textbf{Não funcional (RNF)} & COMO / QUÃO BEM (qualidade, restrição) & ``saldo em tempo real'', desempenho, segurança, usabilidade \\
\bottomrule
\end{tabular}
\end{center}

\textbf{Regra prática:} ``em tempo real'', ``seguro'', ``rápido'',
``disponível'' $\to$ sempre \textbf{não funcional}.

\textbf{Processo:} elicitação $\to$ análise $\to$ especificação $\to$
validação $\to$ gerenciamento. A \textbf{elicitação} (levantamento) usa:
entrevista, questionário, \textbf{brainstorming}, workshop, observação,
prototipação, análise de documentos.

\begin{pegadinha}
A FGV afirma que ``brainstorming é técnica inadequada, use só entrevista
formal'' --- isso é \textbf{falso}; brainstorming é técnica legítima de
elicitação. Outro distrator comum: dizer que o requisito surge \emph{depois}
da implementação (a ordem certa é elicitar antes de codificar).
\end{pegadinha}

\begin{jacaiu}
A FGV também pede para \textbf{contar} quantos RF e quantos RNF há num
enunciado longo, e para classificar o tipo de RNF (usabilidade, desempenho,
produto, organizacional, externo). Ao ler, marque cada verbo de ação (RF) e
cada ``deve ser rápido/seguro/exportável'' (RNF).
\end{jacaiu}

\section{Metodologias ágeis}

\begin{conceito}[Scrum --- é framework, não metodologia]
Trabalho organizado em \textbf{Sprints} de tamanho fixo (time-boxed).

\textbf{Papéis/accountabilities:} \textbf{Product Owner} (valor e Product
Backlog), \textbf{Scrum Master} (remove impedimentos, serve ao time, \emph{não
manda}), \textbf{Developers}.

\textbf{Eventos:} Sprint, Sprint Planning, Daily Scrum, Sprint Review, Sprint
Retrospective.

\textbf{Artefatos:} Product Backlog, Sprint Backlog, Increment.

\begin{itemize}
\item \textbf{Sprint Review = produto} (inspeciona o incremento com
stakeholders).
\item \textbf{Sprint Retrospective = processo} (o que melhorar no jeito de
trabalhar do time).
\item No Daily, com risco à Sprint, o Scrum Master \textbf{facilita a
solução do time} --- não redistribui tarefas sozinho, nem assume a tarefa do
desenvolvedor.
\end{itemize}
\end{conceito}

\begin{center}
\begin{tabular}{@{}p{2.4cm}p{9.8cm}@{}}
\toprule
\textbf{Framework} & \textbf{Ideia central} \\
\midrule
Kanban & Fluxo contínuo, \textbf{sem sprints fechadas}; limita WIP (trabalho em progresso); entrega conforme conclui. \\
XP (Extreme Programming) & Práticas de engenharia: programação em par, TDD, integração contínua, refatoração, cliente presente. \\
Lean & Eliminar desperdício, otimizar o fluxo de valor. \\
Ágil híbrido & Combina práticas ágeis com métodos tradicionais. \\
\bottomrule
\end{tabular}
\end{center}

\begin{pegadinha}
Kanban \textbf{não} tem sprint --- o distrator clássico é ``Kanban organiza o
trabalho em sprints''. Scrum \textbf{tem} time-box; Kanban é fluxo contínuo.
\end{pegadinha}

\begin{jacaiu}
Sprint Planning prioriza pela \textbf{capacidade real da equipe}, nunca por
pressão externa. Mudança tardia de escopo: trata-se com o PO, ajustando o
backlog --- nunca ``recusar por comprometer o plano'' nem ``suspender o
projeto''.
\end{jacaiu}

\section{Testes}

\begin{center}
\begin{tabular}{@{}p{3cm}p{9.2cm}@{}}
\toprule
\textbf{Nível/tipo} & \textbf{O que verifica} \\
\midrule
Unitário & menor unidade isolada (função/método) \\
Integração & interação entre módulos/componentes \\
Sistema & o sistema completo \\
Aceitação & atende ao cliente/requisito \\
Usabilidade & experiência/interface intuitiva \\
Regressão & mudança não quebrou o que já funcionava \\
\bottomrule
\end{tabular}
\end{center}

\begin{conceito}[TDD e BDD]
\textbf{TDD (Test-Driven Development):} escreve o \textbf{teste antes} do
código (ciclo red $\to$ green $\to$ refactor). Foco em design e cobertura.

\textbf{BDD (Behavior-Driven Development):} evolução do TDD, especifica
comportamento em linguagem acessível ao negócio (Gherkin: Dado-Quando-Então).

\textbf{Caixa-preta} (sem ver o código, só entrada/saída) $\times$
\textbf{caixa-branca} (baseado na estrutura interna/código).

\textbf{RPA (Robotic Process Automation):} automatiza tarefas repetitivas de
interface (robôs de software) --- \textbf{não é teste}.
\end{conceito}

\section{Mensuração: Ponto de Função $\times$ Story Points}

\begin{center}
\begin{tabular}{@{}p{2.8cm}p{4.7cm}p{4.7cm}@{}}
\toprule
& \textbf{Ponto de Função (APF)} & \textbf{Story Points} \\
\midrule
Natureza & \textbf{objetiva}, padronizada (IFPUG) & \textbf{relativa/subjetiva}, do time \\
Independe do time? & sim & não (calibrado por time) \\
Bom para & contrato de escopo fechado, comparação entre projetos & estimativa ágil interna \\
Baseado em & funções do sistema (EE, SE, CE, ALI, AIE) & esforço/complexidade percebidos \\
\bottomrule
\end{tabular}
\end{center}

\begin{pegadinha}
A FGV troca os dois: diz que Story Point é objetivo/padronizável entre times,
ou que serve a contrato de escopo fechado --- é o contrário. Ponto de Função é
o objetivo e comparável; Story Points é o relativo do time.
\end{pegadinha}

\begin{comosair}[Classificando telas em APF]
\begin{itemize}
\item Tela que \textbf{ENTRA} e grava dado $\to$ \textbf{EE} (Entrada
Externa).
\item Tela que só \textbf{mostra dado processado/calculado} $\to$
\textbf{SE} (Saída Externa).
\item Tela que só \textbf{consulta e exibe sem cálculo} $\to$ \textbf{CE}
(Consulta Externa).
\item Grupo lógico de dados \textbf{mantido dentro} do sistema $\to$
\textbf{ALI} (Arquivo Lógico Interno).
\item Referência externa só \textbf{lida} (não mantida) $\to$ \textbf{AIE}
(Arquivo de Interface Externa).
\end{itemize}
\end{comosair}

\section{DevOps e entrega}

\begin{itemize}
\item \textbf{CI (Integração Contínua):} integra e testa o código com
frequência (a cada commit).
\item \textbf{CD (Entrega/Implantação Contínua):} leva a nova versão à
produção rapidamente, com mínima interrupção. ``Fornecer rapidamente nova
versão em produção'' = \textbf{CD}, não CI.
\item \textbf{DevOps} une desenvolvimento e operações; \textbf{DevSecOps}
injeta segurança no pipeline.
\end{itemize}

\section{Design de software}

\begin{itemize}
\item \textbf{Design de alto nível} = estrutura geral, módulos e sua
interação (próximo da arquitetura). \textbf{Design de baixo nível} = detalhe
de funções, métodos, algoritmos.
\item \textbf{Arquitetura} = decisões amplas e estruturais; \textbf{design} =
decisões detalhadas. Não é verdade que ``arquitetura só serve para projeto
grande''.
\item Padrões de projeto e UML têm capítulo próprio (Capítulo
\ref{cap:padroes-uml}).
\end{itemize}

\begin{jacaiu}
Cascata $\times$ incremental; Sprint Review $\times$ Retrospective; requisito
funcional $\times$ não funcional (cenário do ``saldo em tempo real'');
brainstorming como elicitação válida; papel do Scrum Master no Daily e no
Sprint Planning; metodologia ágil para mudança frequente (XP/Scrum/Kanban); CD
no DevOps; Ponto de Função $\times$ Story Points; design alto $\times$ baixo
nível; níveis de teste + TDD; ágil híbrido. Rode \texttt{./quiz.py
eng-software}.
\end{jacaiu}

\begin{pegadinha}[Resumo dos pares que a FGV inverte neste bloco]
Review$\leftrightarrow$Retro, funcional$\leftrightarrow$não funcional,
PF$\leftrightarrow$Story Points, cascata$\leftrightarrow$incremental,
TDD$\leftrightarrow$BDD, alto$\leftrightarrow$baixo nível. Absolutos
(``sempre'', ``exclusivamente'', ``apenas'', ``impossível'') quase sempre
marcam o distrator. E cuidado com distorcer o papel do Scrum Master: ele
facilita e serve, nunca comanda nem executa a tarefa do desenvolvedor.
\end{pegadinha}

\section{Alta probabilidade / pesquisa extra}

\begin{itemize}
\item \textbf{Scrum Guide 2020:} usa ``accountabilities'' (não ``papéis'');
removeu a figura do ``time de desenvolvimento'' como subgrupo --- hoje é um
\textbf{Scrum Team} único com Developers. Sprint Goal, Product Goal,
Definition of Done.
\item \textbf{Clean Code / SonarQube} (no edital): análise \textbf{estática}
de código (sem executar) para dívida técnica, code smells, cobertura.
\item \textbf{MVP, design thinking, UX} aparecem no edital do Perfil 1 mas
permeiam o Perfil 3; MVP = menor versão que entrega valor e valida uma
hipótese.
\end{itemize}
